\subsection{四単純位数4のpaired rank一型の全排除}
\label{subsec:cubic-four-simple-order4-rankA}
\begin{theorem}[二つ目のrank型の完成]
\label{thm:cubic-four-simple-order4-rankA-empty}
スペクトル $(1,-1,i,-i)$ の必要rank型 $(2,2,1,1)$ も全てH3を与えない。
従ってこのスペクトルに残る必要rank型は $(2,0,2,2)$ だけである。
\end{theorem}
\begin{proof}
前節の記号を使う。$M$ はrank一なので、非零列 $P$ と非零行 $Q^t$ により
$M=PQ^t$ と書ける。
$P$ の基底は $(xb,yc,bx,cy)$、$Q$ の基底は $(xc,yb,by,cx)$ である。
いずれも四つのtensorは相異なる行と列に位置するので、
一本だけ非零成分を持つ $P$ 又は $Q$ は純積関係となる。
整域性から両supportの大きさは2以上である。
元巡回式は正確に
\begin{gather*}
 p_2q_3=-p_0q_0,\qquad p_2q_2=ip_0q_1,\\
 p_3q_3=-ip_1q_0,\qquad p_3q_2=p_1q_1
\end{gather*}
を与える。

この四式の非零supportを調べると、可能なものは次の十組だけである。
$02$ はsupport $\{0,2\}$、$0123$ は全supportの略記とする。
\[
\begin{array}{c|c}
 \operatorname{supp}P&\operatorname{supp}Q\\\hline
 02&03,\ 12,\ 0123\\
 13&03,\ 12,\ 0123\\
 0123&03,\ 12\\
 01&23\\
 23&01
\end{array}
\]
実際、各式で二つの積の一方だけが非零となるsupportは不可能である。
大きさ2以上のsupportの全組についてこの条件を適用すると、
上の十組と両方全supportの場合だけが残る。
最後の場合は最初と第三式から $p_2/p_3=-ip_0/p_1$、
第二と第四式から $p_2/p_3=ip_0/p_1$ となり矛盾する。
これはsupportに関する有限の通常の議論である。
公開再現器は全組を列挙し、非零座標をsaturateした式でも同じ十組を再現する。

$b,c$ を交換し、原始四乗根を $i'=-i$ と読み替えると、
元のスペクトルと巡回式は同じ形を保つ。
この操作は $P02$ と $Q03$、$P13$ と $Q12$、
$P01,Q23$ と $P23,Q01$ を交換する。
従って以下の三chartを調べれば全被覆となる。
これは対象のintrinsic同型による軌道同一視を主張したものではない。

\paragraph{$P02$ chart．}
外積のscalar自由度を使って
\[
 P=(1,0,p,0)^t,\quad
 Q=(u,v,iv/p,-u/p)^t,\qquad p\ne0
\]
と書ける。$H_2$ の $c^2$ 列は $a_{14}b^2$ なので、整域性から $a_{14}=0$。
$H_2$ のrank2と整域性を合わせると $a_0v\ne0$。
実際、$v=0$ なら純 $xy,yx$ 関係が生じ、
$a_0=0,v\ne0$ なら純 $b^2$ 関係が生じ、両方零ならrankが零となる。
このchartの重み0 middleは
\[
 H_0=\begin{pmatrix}
 0&a_0&u/p&-ipu\\a_0&a_{10}&0&0\\
 u/p&0&a_{15}&0\\-ipu&0&0&-ia_{15}
 \end{pmatrix}.
\]
$u=0$ なら第一列は非零純 $y^2$ 関係である。
$u\ne0$ なら左上 $2\times2$ blockは可逆なので、rank2はそのSchur補の消滅と同値。
その非対角成分は $-ia_{10}u^2/a_0^2$ だから $a_{10}=0$、
対角成分から $a_{15}=0$ となる。
すると第二列は非零純 $x^2$ 関係となり、やはり整域性に反する。

\paragraph{$P13$ chart．}
同様に
\[
 P=(0,1,0,p)^t,\quad
 Q=(u,v,v/p,-iu/p)^t,\qquad p\ne0
\]
と書ける。$H_0$ の第一列は $a_0y^2$ なので $a_0=0$。
$H_2$ の $b^2$ 列は $-a_{14}c^2$ なので $a_{14}=0$。
この二つを消すと $H_2$ の第一列は $-iu c^2/p$ である。
$u\ne0$ なら純平方関係、$u=0$ なら $H_2=0$ となり、いずれも不可能。

\paragraph{$P01,Q23$ chart．}
外積と非零topのscalar変更で
\[
 P=(1,p,0,0)^t,\quad Q=(0,0,1,q)^t,\qquad pq\ne0
\]
と書ける。$H_2$ のPfaffianは $-a_0a_{14}-pq$ なので
$a_0a_{14}=-pq\ne0$ が必要である。
一方、$H_0$ は
\[
 \begin{pmatrix}
 0&a_0&-q&0\\a_0&a_{10}&0&p\\
 -q&0&a_{15}&-ia_{14}\\0&p&-ia_{14}&-ia_{15}
 \end{pmatrix}.
\]
左上blockが可逆であるから、rank2はSchur補の消滅を強制する。
右下成分から $a_{15}=0$、左上成分から $a_{10}=0$。
非対角成分は $-ia_{14}+pq/a_0$ であり、Pfaffian条件を代入すると
\[
 (1+i)pq/a_0\ne0
\]
となって矛盾する。
以上で六chartの全てが消え、paired rank一型は空である。
\end{proof}

公開再現器は \texttt{verify\_cubic\_four\_simple\_order4\_rankA.py}。
全support、巡回四式、対称操作、元16tensorからの三chartの復元、
純積列、PfaffianとSchur恒等式を標数0で検算する。
正常核の一般点計算や有限標本に依存せず、必要rankと整域性だけで全値を排除した。
残る $(2,0,2,2)$ と全四単純・全CY・相対分類は別の未完成問題である。
